\section{Screening Certificate Details}\label{sec:screenec53}
The screening optimizer evaluates the anchor envelope once per candidate. The verification implementation is deliberately different. For every declared deletion it enumerates all remaining catalog predecessors and all possible successors, including the terminal successor. For each history admitting the deleted command, it computes the difference between the two canonical responses at the cap and every eligible knot in the local interval. With quadratic primitives, this difference is linear on common interpolation pieces, convex where one response interpolates and the other uses service, and linear where both have quadratic service tails with the same curvature. Its maximum is therefore attained at these endpoints. The checker sums the individual maxima and maximizes over neighboring contexts. Theorem~\ref{thm:screen53} proves that this result equals the anchor envelope. Verification uses rational arithmetic and imports neither the screening optimizer nor its bound formula.

The implementation uses a descending deletion order, although the theorem permits any order. A record names the candidate, charge, bound, strictness flag, and the anchor--terminal maximizing context. A missing, repeated, or potentially anchoring deletion fails. The resulting catalog and its charges must match the catalog and charges in the supplied resource certificate exactly. Histories, promise, curvature parameters, and requested tolerance must be unchanged. The reduced certificate uses the minimum of the original budget and the retained catalog size, which describes the same cardinality restriction. The retained lower policy then satisfies the full original constraints.

Computing all screening bounds requires $O(kN)$ arithmetic operations and $O(N)$ recorded scalar witnesses. The reference screening checker uses at most $O(kN^3)$ operations over all candidate deletions because it explicitly enumerates neighboring contexts; the independent resource checker has its own table-verification cost. This deliberately redundant verification path is polynomial but is not advertised as the fastest available implementation. A linear-time bound checker could instead evaluate the proved envelope directly. Retaining the original model also places a lower bound on certificate storage.

An expected-instance digest supplied outside the certificate prevents accepting a valid certificate for a different mathematical instance. The embedded digest detects accidental inconsistency but is not an authentication mechanism. Verification still checks the mathematical contents. Python's optimization flags are rejected by the inherited resource checker; disabling assertions is not a supported verification mode.

The first development run completed the random regression but stopped at an overly restrictive test expectation that the four base candidates would all remain. The actual rule additionally deleted the globally ineligible, zero-charge endpoint one. The correction changed only that expected retained set, not the model, screening criterion, protocol parameters, or measured inputs. The original diagnostic is preserved with the execution record. All twelve cases were then rerun. This distinction avoids treating a corrected test expectation as a change to the experiment.

The complete controlled records contain exact input digests, exhaustive values, both global intervals, selected books, class counts, state counts, compressed and uncompressed bytes, optimization times, screening times, and verification times. The original full certificate and the transferred reduced certificate are stored for each case. Retained earlier experiments and unresolved historical requests are not overwritten by these new records.
